\documentclass[10pt,a4paper]{amsart}
\usepackage[margin=19mm]{geometry}
\usepackage{amsmath,amssymb,mathtools}
\usepackage[scr=rsfs,cal=euler]{mathalfa}
\usepackage{xfrac}
\usepackage[unicode,hidelinks,bookmarks=false]{hyperref}
\newcommand{\ie}{i.e.\ }
\newcommand{\cf}{cf.\ }
\newcommand{\resp}{respectively\ }
\newcommand{\Subsection}{\subsection}
\providecommand{\nobreakdash}{\nobreak\hbox{-}}
\setlength{\emergencystretch}{2em}
\multlinegap0pt
\usepackage{bm}
\usepackage{bbm}
\usepackage{dsfont}
\usepackage{xspace}
\usepackage{cancel}
\usepackage{mathtools}
\usepackage{amsthm}
\makeatletter
\@ifundefined{theorem}{\newtheorem{theorem}{Theorem}}{}
\@ifundefined{lemma}{\newtheorem{lemma}[theorem]{Lemma}}{}
\makeatother
\DeclareMathOperator{\Sq}{Sq}
\newcommand{\Z}{\mathbb{Z}}
\newcommand{\supp}{\operatorname{supp}}
\title[Small covers as pullbacks from the simplex]{Small covers as pullbacks from the simplex}
\author{Suyoung Choi, Hyeontae Jang, Younghan Yoon}
\date{Source: arXiv:2604.13437v1 (CC BY 4.0)}
\begin{document}
\maketitle

\noindent\emph{Source and licence.} Small covers as pullbacks from the simplex. Version arXiv:2604.13437v1 (CC BY 4.0), distributed under the Creative Commons Attribution 4.0 International licence, as recorded in the arXiv licence metadata for this exact version and shown on the abstract page https://arxiv.org/abs/2604.13437v1. Full rights notice: licenses/arxiv-2604.13437-CC-BY-4.0.txt.\par\smallskip

\noindent\textit{Editorial scope.} Counted unit: the Theorem whose source label is \texttt{thm:Sq1\_H2\_characterizes\_simplex\_pullback} in the article (source0.tex lines 471--479, source label \texttt{thm:Sq1\_H2\_characterizes\_simplex\_pullback}), delivered together with the article's own complete original proof of it (source0.tex lines 481--552, one \texttt{proof} environment of 72 lines). No mathematics is written, repaired or re-derived: statement and proof are the article's own text line for line. Required context retained uncounted: lem:square\_generator\_RT (lines 421--430). Every definition, display and label of those lines is retained unchanged. Omitted from this package and not redistributed: 1--420, 431--470, 480--480, 553--934. Nothing inside a retained excerpt is omitted or altered: every retained body is delivered exactly as the article's lines are written, so no omission receipt of any kind accompanies this package. Citations: the retained excerpts carry no citation command at all, so no bibliography entry of the article is transcribed and no reference is redistributed. Reference adaptation: none was needed and none was made; every reference inside the delivered unit resolves inside it, so no unresolved reference or citation is produced. Printed slips kept literal: no printed word, symbol, hypothesis or number of the counted statement or of its proof was altered. Layout and class: the article's own class is \texttt{amsart} with packages including amsmath, amssymb, amsthm, mathtools, color, amscd, hyperref, array, booktabs; the wrapper is the lane's pinned \texttt{amsart} editorial wrapper at 10pt on A4, which restates the theorem-like environments the retained text needs and adds the article's own macro definitions that the retained text needs (\texttt{\textbackslash Sq}, \texttt{\textbackslash Z}, \texttt{\textbackslash supp}), with extra wrapper packages (bm, bbm, dsfont, xspace, cancel, mathtools). The delivered numbering is the unit's own. Assets: no retained span contains a figure environment, a graphics-inclusion command, a PDF-page inclusion, a table or a TikZ picture; no image file is copied into this package, the package declares no asset dependency and no image file is redistributed with it. Licence: version arXiv:2604.13437v1 of this article is distributed under the Creative Commons Attribution 4.0 International licence, as recorded in the arXiv licence metadata for this exact version and shown on the abstract page https://arxiv.org/abs/2604.13437v1. A full rights notice is delivered at licenses/arxiv-2604.13437-CC-BY-4.0.txt. Characters: every retained excerpt is pure ASCII, so the delivered engine, XeLaTeX with its default Latin Modern text font, reports no missing character.
\begin{lemma}\label{lem:square_generator_RT}    
    For every generator $v_j\in H^1(M;\Z_2)$,
    $$
    v_j^2=\tau v_j.
    $$
    In particular, for a homogeneous element $x \in  H^q(M;\Z_2)$, we have
    $$
        \Sq(x)=(1+ \tau)^q x.
    $$
\end{lemma}

\begin{theorem}\label{thm:Sq1_H2_characterizes_simplex_pullback}
    Let~$M$ be a small cover.
    Then the following are equivalent.
    \begin{enumerate}
        \item\label{eq:FT1} $M$ is a pullback from the simplex;
        \item\label{eq:FT2} $\Sq^1$ vanishes on $H^{2k}(M;\Z_2)$ for all $k\geq 0$; and
        \item\label{eq:FT3} $\Sq^1$ vanishes on $H^2(M;\Z_2)$.
    \end{enumerate}
\end{theorem}

\begin{proof}
    The implication \eqref{eq:FT1} $\Rightarrow$ \eqref{eq:FT2} follows from Lemma~\ref{lem:square_generator_RT}, since for $x \in H^q(M;\Z_2)$
    $$
    \Sq^1(x) = \binom{q}{1}\tau x,
    $$
    which vanishes whenever $q$ is even.
    The implication \eqref{eq:FT2}$\Rightarrow$\eqref{eq:FT3} is obvious.
    
    Let us show the implication \eqref{eq:FT3} $\Rightarrow$ \eqref{eq:FT1}.
    Assume that $M(K,\lambda)$ is not a pullback from the simplex. 
    Then, there exist a facet $\sigma=\{u_1,\dots,u_n\}$ and an index $i\in[n]$ such that, after a change of basis of~$\Z_2^n$, we have
    $$
        \lambda(u_1)=e_1,\dots,\lambda(u_n)=e_n,
    $$ and $\lambda(p_i (\sigma))=\sum_{j\in S_i^\sigma}e_j$ for a subset $S_i^\sigma \subset[n]$ satisfying
    $$
        S_i^\sigma \neq \{i\} \quad \text{ and }\quad S_i^\sigma \neq [n].
    $$
    
    Set $p:=p_i (\sigma)$.
    Since $(\sigma\setminus\{u_i\})\cup\{p\}$ is a facet, the vector $\lambda(p)$ is not contained in the span of $\{e_j\mid j\neq i\}$, so necessarily $i\in S_i^\sigma$.
    Choose
    $$
        s\in S_i^\sigma \setminus\{i\} \quad \text{ and } \quad t \in [n]\setminus S_i^\sigma.
    $$
    We claim that
    $$
        \Sq^1(v_{u_s} v_{u_t})= v_{u_s} v_{u_t}(v_{u_s} + v_{u_t})\neq 0.
    $$
    
    For a vertex $q \in V(K)$, let $\supp_\sigma(q)\subset[n]$ be the support of $\lambda(q)$ with respect to the basis
    $$
        \lambda(u_1)=e_1,\dots,\lambda(u_n)=e_n,
    $$
    that is,
    $$
        \lambda(q)=\sum_{j\in \supp_\sigma(q)} e_j.
    $$
    
    For each $r\in[n]$, the $r$th row relation in $J_\lambda$ gives
    $$
        v_{u_r} = \sum_{\substack{q\notin \sigma\\ r\in \supp_\sigma(q)}} v_q.
    $$
    Hence
    $$
    v_{u_s}+v_{u_t} = \sum_{\substack{q\notin \sigma\\ |\supp_\sigma(q)\cap\{s,t\}|=1}} v_q.
    $$
    Since $s\in S_i^\sigma$ and $t\notin S_i^\sigma$, the vertex $p_i(\sigma)$ appears in this sum.
    
    Now let
    $$
        R':=\prod_{r\in[n]\setminus\{i,s,t\}} v_{u_r} \quad \text{ and } \quad R := R' v_{u_s} v_{u_t}.
    $$
    Then
    $$
        R' \Sq^1(v_{u_s}v_{u_t}) = R(v_{u_s}+v_{u_t}) = \sum_{\substack{q\notin \sigma\\ |\supp_\sigma(q)\cap\{s,t\}|=1}} Rv_q.
    $$
    A term $Rv_q$ is nonzero if and only if
    $$
        \bigl(\sigma\setminus\{u_i\}\bigr)\cup\{q\}
    $$
    is a facet.
    Since $K$ is a simplicial sphere, each codimension one face is contained in exactly two facets. 
    Hence the only nonzero term in the above sum is the one corresponding to $q=p$, and so
    $$
        R'  \Sq^1(v_{u_s}v_{u_t}) = R v_{p}\neq 0.
    $$
    Therefore
    $$
        \Sq^1(v_{u_s}v_{u_t})\neq 0
    $$
    as desired.
\end{proof}

\end{document}
